\section{Neighboring Methods and Assumptions}\label{sec:min-related}

C-TreePO sits between three literatures: design-based supervised learning
in social-science measurement, mergeable algebraic aggregation in
data-systems, and preference optimization in language-model alignment.
Appendix~\ref{app:min-related} records the longer scope notes;
Section~\ref{sec:min-related} positions the framework against the
neighboring methods that share pieces of the workflow but do not provide
the same local-to-global certificate.

\paragraph{Design-Based Supervised Learning.}
\citet{EgamiEtAl2024} introduce the design-based-supervised-learning (DSL)
framework for social-science applications of LLMs. They show that high
prediction quality on surrogate labels is not enough when those labels
enter downstream inference: errors can be correlated with the estimand,
and a logged sampling design over a small gold set is what makes the
imputation provably valid. C-TreePO's audit is the same idea applied one
layer earlier in the pipeline, at the compression step rather than at
final-label imputation. The local laws name what has to be preserved by
each compression call, and the IPW estimator over realized tree nodes
turns realized compression into the kind of design-based estimand DSL
already trusts at the document level. Related lineages on
imperfect-surrogate inference, measurement validity, and structural-topic
modeling \citep{RobertsStewartTingley2014} share the construct-validity
framing the audit operationalizes.

\paragraph{Long-Context Methods.}
RAG retrieves relevant chunks but provides no guarantee about what is
\emph{lost} in retrieval; for non-separable targets, retrieval may miss
interaction-bearing information that only becomes decision-relevant when
chunks are combined. Sliding-window methods, recursive summarization, and
extended-context models similarly lack auditable certificates. C-TreePO
processes all chunks through a tree and uses C3-style audits to certify
whether cross-chunk interactions are preserved in the realized pipeline,
providing a certification layer that complements these practical
approaches.

\paragraph{Tree-Based Text Composition and Compression.}
Tree-structured inference has been used to compose and compress
natural-language descriptions, with classical work on sentence compression
and fusion providing early structured formulations
\citep{KnightMarcu2000SentenceCompression,BarzilayMcKeown2005SentenceFusion,ClarkeLapata2008ILPSentenceCompression}.
\citet{KuznetsovaEtAl2014TreeTalk} harvest syntactic tree fragments from
web captions and use ILP- and CKY-style structured inference to compose
new image descriptions and generalize captions via parse-tree pruning.
The closest conceptual resonance is that this line already treats
composition and compression together in one tree-based workflow. C-TreePO
differs by making the compressed node a task-relative state and by adding
an auditable local-to-global certificate for oracle preservation and
preference-objective equivalence.

\paragraph{Sufficient Statistics, Mergeable Summaries, and Sketch-Flip-Merge.}
The local-law machinery sits at the end of a four-generation lineage
of approximate-sufficiency results. Fisher and Blackwell define
sufficient statistics and the sufficiency ordering of experiments
\citep{Fisher1922,Blackwell1953}: a statistic is sufficient for a
parameter when the conditional distribution of the data given the
statistic does not depend on the parameter. Mergeable summaries
\citep{Agarwal2013MergeableTODS,FlajoletEtAl2007,CormodeMuthukrishnan2005,GreenwaldKhanna2001}
add an associative merge axiom and a sublinear-state requirement;
Count-Min, HyperLogLog, Greenwald-Khanna quantiles, and Bloom filters
each pair an exact sufficient statistic with a fixed query function,
and correctness is proved once algebraically. Sketch-Flip-Merge
\citep{HehirTingCormode2023SFM} extends the mergeable abstraction to
local differential privacy: sufficiency must survive a local
randomizer between leaf and merge, and only matched
randomizer-summarizer pairs keep the merged channel in a parametric
family the cardinality estimator can invert.

C-TreePO inherits the encode-merge-readout discipline. It replaces
hand-designed sketch state with a learned summarizer, replaces the
closed-form sufficiency proof with an audit that estimates residual
local-law violations, and replaces opaque register state with text
summaries that are inspectable in the same units as the leaf.
Theorem~\ref{thm:sketch-equivalence} states the equivalence between
the local laws and the mergeable-summary algebra; the proof is in
Appendix~\ref{app:min-algebra} and is verified in Lean
(\texttt{SketchRecovery.lean}). The empirical HLL parity result is in
Appendix~\ref{app:min-hll}, and Appendix~\ref{app:min-related} records
the longer sufficient-statistics scope notes.

\paragraph{Preference Learning and RLHF.}
Direct Preference Optimization fine-tunes a language model from pairwise
human preferences by reparameterizing the Bradley--Terry likelihood so
that the policy itself serves as an implicit reward model
\citep{BradleyTerry1952,RafailovEtAl2023DPO}. The results above apply to
DPO and GRPO variants when the preference signal depends on the input only
through the preserved oracle value. When readability, fluency, or style
preferences fall outside $\fstar$, they must be modeled by a richer oracle
or handled by separate alignment channels.

\paragraph{Assumption-Centered Guarantees.}
Recent machine-learning theory often separates structural invariance
claims from behavioral assumptions on annotator noise and loss geometry.
The analysis here follows that pattern but adds an explicit
representational layer: neural-operator approximation and transfer moduli
propose local-law budgets, C1--C3 and context compatibility drive the
preservation/equivalence results, and quantitative transport to preference
gaps uses explicit method-level Lipschitz and integrability conditions.
Keeping this split explicit makes guarantees easier to audit and
stress-test in applied settings.
